% Options for packages loaded elsewhere
\PassOptionsToPackage{unicode}{hyperref}
\PassOptionsToPackage{hyphens}{url}
\PassOptionsToPackage{dvipsnames,svgnames,x11names}{xcolor}
\documentclass[
  a4paper,
]{article}
\usepackage{xcolor}
\usepackage[margin=20mm]{geometry}
\usepackage{amsmath,amssymb}
\setcounter{secnumdepth}{-\maxdimen} % remove section numbering
\usepackage{iftex}
\ifPDFTeX
  \usepackage[T1]{fontenc}
  \usepackage[utf8]{inputenc}
  \usepackage{textcomp} % provide euro and other symbols
\else % if luatex or xetex
  \usepackage{unicode-math} % this also loads fontspec
  \defaultfontfeatures{Scale=MatchLowercase}
  \defaultfontfeatures[\rmfamily]{Ligatures=TeX,Scale=1}
\fi
\usepackage{lmodern}
\ifPDFTeX\else
  % xetex/luatex font selection
  \setmainfont[]{DejaVu Serif}
  \setsansfont[]{DejaVu Sans}
  \setmonofont[]{DejaVu Sans Mono}
\fi
% Use upquote if available, for straight quotes in verbatim environments
\IfFileExists{upquote.sty}{\usepackage{upquote}}{}
\IfFileExists{microtype.sty}{% use microtype if available
  \usepackage[]{microtype}
  \UseMicrotypeSet[protrusion]{basicmath} % disable protrusion for tt fonts
}{}
\makeatletter
\@ifundefined{KOMAClassName}{% if non-KOMA class
  \IfFileExists{parskip.sty}{%
    \usepackage{parskip}
  }{% else
    \setlength{\parindent}{0pt}
    \setlength{\parskip}{6pt plus 2pt minus 1pt}}
}{% if KOMA class
  \KOMAoptions{parskip=half}}
\makeatother
\usepackage{longtable,booktabs,array}
\newcounter{none} % for unnumbered tables
\usepackage{calc} % for calculating minipage widths
% Correct order of tables after \paragraph or \subparagraph
\usepackage{etoolbox}
\makeatletter
\patchcmd\longtable{\par}{\if@noskipsec\mbox{}\fi\par}{}{}
\makeatother
% Allow footnotes in longtable head/foot
\IfFileExists{footnotehyper.sty}{\usepackage{footnotehyper}}{\usepackage{footnote}}
\makesavenoteenv{longtable}
\usepackage{graphicx}
\makeatletter
\newsavebox\pandoc@box
\newcommand*\pandocbounded[1]{% scales image to fit in text height/width
  \sbox\pandoc@box{#1}%
  \Gscale@div\@tempa{\textheight}{\dimexpr\ht\pandoc@box+\dp\pandoc@box\relax}%
  \Gscale@div\@tempb{\linewidth}{\wd\pandoc@box}%
  \ifdim\@tempb\p@<\@tempa\p@\let\@tempa\@tempb\fi% select the smaller of both
  \ifdim\@tempa\p@<\p@\scalebox{\@tempa}{\usebox\pandoc@box}%
  \else\usebox{\pandoc@box}%
  \fi%
}
% Set default figure placement to htbp
\def\fps@figure{htbp}
\makeatother
\setlength{\emergencystretch}{3em} % prevent overfull lines
\providecommand{\tightlist}{%
  \setlength{\itemsep}{0pt}\setlength{\parskip}{0pt}}
\usepackage{luatexja}
\usepackage{xurl}
\usepackage{float}
\usepackage{bookmark}
\IfFileExists{xurl.sty}{\usepackage{xurl}
\usepackage{float}}{} % add URL line breaks if available
\urlstyle{same}
\hypersetup{
  pdftitle={Emergent Quantum Structure from the 3D Triangle-Loop Product: Exact Discretization{,} Spontaneous Single Complex Structure{,} Born Rule{,} and Discrete Uncertainty Saturation (Fourth Rapid Communication)},
  pdfauthor={Noriaki Kihara (WF System Co., Ltd.) ORCID 0009-0004-6753-4020},
  colorlinks=true,
  linkcolor={Maroon},
  filecolor={Maroon},
  citecolor={Blue},
  urlcolor={Blue},
  pdfcreator={LaTeX via pandoc}}

\title{Emergent Quantum Structure from the 3D Triangle-Loop Product:
Exact Discretization, Spontaneous Single Complex Structure, Born Rule,
and Discrete Uncertainty Saturation (Fourth Rapid Communication)}
\usepackage{etoolbox}
\makeatletter
\providecommand{\subtitle}[1]{% add subtitle to \maketitle
  \apptocmd{\@title}{\par {\large #1 \par}}{}{}
}
\makeatother
\subtitle{― Analytical Derivation of Complex Hilbert Structure,
Probability-Squared Law, and Minimum Resolution from a Real 3-Tensor
Without Quantum Postulates ―}
\author{Noriaki Kihara (WF System Co., Ltd.) ORCID 0009-0004-6753-4020}
\date{November 28, 2026}

\begin{document}
\maketitle

\textbf{Author:} Noriaki Kihara (WF System Co., Ltd.)\\
\textbf{ORCID:} 0009-0004-6753-4020\\
\textbf{Date:} November 28, 2026\\
\textbf{Status:} Fourth Rapid Communication following Rapid Reports 1
(DOI: 10.5281/zenodo.24616931), 2 (DOI: 10.5281/zenodo.24618400), and 3
(DOI: 10.5281/zenodo.24619357). Establishes four analytical theorems
with machine-precision numerical verification showing that the core
pillars of quantum mechanics emerge directly from a purely real 3-tensor
iteration without quantum postulates.

\begin{center}\rule{0.5\linewidth}{0.5pt}\end{center}

\subsection{Abstract}\label{abstract}

We investigate the background-free, purely real 3-tensor triangle-loop
product \(\mathrm{Tri}(T)_{ijk} = \sum_{p,q,r} T_{ipq} T_{qjr} T_{rpk}\)
without assuming spacetime, an action, complex amplitudes, or
probabilistic postulates. Synthesizing the \(N = 64\)
(\(O_8 \otimes O_8\)) results established in Rapid Communications 1--3,
we prove and numerically verify four foundational theorems showing that
quantum structure emerges intrinsically from the algebra itself:

\begin{enumerate}
\def\labelenumi{\arabic{enumi}.}
\tightlist
\item
  \textbf{Theorem Q1 (Exact Discretization)}: The cubic phase map
  \(\theta \mapsto 3\theta\) combined with the null/chiral balance
  condition \(\sum z_n^2 = 0\) quantizes periodic phases to rational
  angles \(2\pi\mathbb{Z}/(3^p - 1)\), restricts spatial crystals to
  Hurwitz dimensions (\(1, 2, 4, 8\)), and discretizes internal
  view-angle couplings \(\rho = L_3/L_1\) to exactly 7 values.
\item
  \textbf{Theorem Q2 (Spontaneous Single Complex Structure
  \(\mathbb{C}\))}: Although the system starts from a real array
  \(T \in \mathbb{R}^{64\times 64\times 64}\), the inner sector \(M_8\)
  of Layer 1 (\(H = O_8 \otimes M_8\)) suppresses symmetric anisotropic
  Pauli components (\(\sigma_x, \sigma_z\)) to \(\sim 10^{-16}\),
  spontaneously selecting a pure complex field
  \(\mathbb{C} \cong \{c_0 I + c_J(i\sigma_y)\}\). Furthermore, Layer 2
  (\(\mathbb{C}_{\mathrm{out}} \otimes_{\mathbb{R}} \mathbb{C}_{\mathrm{in}} \cong \mathbb{C}_+ \oplus \mathbb{C}_-\))
  is unstable and collapses into Layer 3,
  \(U_{\text{Layer 3}} = \frac{1}{\sqrt{2}}(u\otimes v + (Ju)\otimes(-Jv))\),
  locking a single global imaginary unit \(J\) while exactly doubling
  \(\|P\|\) (\(|\sigma_{1,2} - 1/\sqrt{2}| = 3.23 \times 10^{-7}\)).
\item
  \textbf{Theorem Q3 (\(O(N)\) Covariance, \(L^2\) Uniqueness, and the
  Born Rule)}: The map \(\mathrm{Tri}\) is strictly \(O(64)\)-covariant
  (relative error \(8.84 \times 10^{-16}\)). Among \(L^p\) norms, only
  \(p = 2\) is invariant under \(O(64)\) rotations (relative change
  \(2.22 \times 10^{-16}\), vs.~\(+351.8\%\) for \(L^1\) and \(-42.0\%\)
  for \(L^4\)). Together with the emergent complex structure
  \(\mathbb{C}\) (Theorem Q2) and dimension \(d \ge 3\), Gleason's
  theorem uniquely selects the Born rule
  \(P(\Pi) = \mathrm{Tr}(\rho\Pi) = |\langle\phi|\psi\rangle|^2\).
\item
  \textbf{Theorem Q4 (\(\mathbb{Z}_2^3\) Discrete Uncertainty Saturation
  and Minimum Resolution)}: On the outer crystal
  \(O_8 \cong \mathbb{Z}_2^3\), any normalized state \(f\) and its
  Walsh--Hadamard spatial harmonic transform \(\hat{f} = H_8 f\) satisfy
  the exact Donoho--Stark bound
  \(|\mathrm{supp}(f)|\cdot|\mathrm{supp}(\hat{f})| \ge 8\) and Shannon
  entropic uncertainty bound
  \(H_{\mathrm{space}}(f) + H_{\mathrm{harmonic}}(\hat{f}) \ge \ln 8 \approx 2.079442\).
  All three bare localization prototypes---nonlocal (\(8\times 1\)),
  quaternion core \(H_4\) (\(4\times 2\)), and 1-site core
  \(\delta_{000}\) (\(1\times 8\))---strictly saturate both bounds
  (excess \(0.00 \times 10^0\)), whereas \(\tau = 30\) dressed states
  (\(L_x\), P5, P4) shift strictly above the bound (\(+0.051\) to
  \(+0.313\)) due to their background clouds. Minimum spatial resolution
  \(\Delta x_{\min} = 1\) site (Planck length \(\ell_P\) prototype), UV
  cutoff \(k_{\max} = 3\), minimum step \(\Delta\tau_{\min} = 1\)
  (Planck time \(t_P\) prototype), and minimum phase-space cell
  \(\ln 8\) (dimensionless \(\hbar\) prototype) arise directly from the
  algebra.
\end{enumerate}

\begin{center}\rule{0.5\linewidth}{0.5pt}\end{center}

\subsection{1. Introduction and Scope}\label{introduction-and-scope}

In Rapid Communications 1--3 {[}1--3{]}, we showed that iterating the
normalized triangle-loop product
\(U^{(\tau+1)} = \mathrm{Tri}(U^{(\tau)})/\|\mathrm{Tri}(U^{(\tau)})\|_F\)
on \(N = 64\) (\(O_8 \otimes O_8\)) selects Hurwitz crystals, undergoes
a universal four-stage doubling-entanglement phase transition (rank
\(1 \to 2 \to 4\)), exhibits \(SO(3)\) spin selection (\(1\)
vs.~\((1/3)^\tau\)) and \(G_2\) (\(14\)) \(\supset SU(3)\) (\(8\)) gauge
symmetries, separates nonlocal gauge waves from localized condensates
via \(\mathbb{Z}_2^3\) spatial harmonics, and supports two-particle
coexistence.

A central foundational question remains: is this merely a classical
nonlinear tensor map, or does it inherently generate the mathematical
axioms of quantum mechanics? Here we demonstrate---via four analytical
theorems and deterministic machine-precision verification
(\texttt{tri\_loop\_n64\_quantum\_foundations\_verify\_v0.py}, zero
randomness)---that discretization, a single complex Hilbert structure,
the Born rule, and discrete uncertainty saturation with a UV cutoff
emerge without external quantum postulates. Throughout, \(\tau\) denotes
the discrete update step, distinct from emergent physical time \(t\).

\begin{center}\rule{0.5\linewidth}{0.5pt}\end{center}

\subsection{\texorpdfstring{2. Theorem Q1: Exact Discretization via
\(\sum z_n^2 = 0\) and Cubic
Self-Consistency}{2. Theorem Q1: Exact Discretization via \textbackslash sum z\_n\^{}2 = 0 and Cubic Self-Consistency}}\label{theorem-q1-exact-discretization-via-sum-z_n2-0-and-cubic-self-consistency}

\textbf{Theorem Q1 (Phase, Coupling, and Crystal Discretization)}:
Because \(\mathrm{Tri}\) is homogeneous of degree 3, phase evolution
acts as the tripling map \(\theta_n \mapsto 3\theta_n \pmod{2\pi}\). Any
self-consistent fixed point or period-\(p\) cycle is therefore
restricted to discrete rational angles:
\[(3^p - 1)\theta_n \equiv 0 \pmod{2\pi} \quad \Longrightarrow \quad \theta_n \in \frac{2\pi}{3^p - 1}\mathbb{Z}.\]
Under the null/chiral condition
\(\sum_n z_n^2 = \sum_n r_n^2 e^{2i\theta_n} = 0\) (the \(I + \sigma_y\)
sector in real representation {[}2{]}), arbitrary continuous phase
distributions cannot balance the quadratic invariant, forcing discrete
polygonal equilibria. Spatially, Hurwitz's theorem restricts
self-consistent division-algebra crystals to dimensions \(1, 2, 4, 8\),
and the internal triangle view-angle couplings \(\rho = L_3/L_1\) of
Layer 1 (\(H = O_8 \otimes M_8\)) quantize into exactly 7 discrete
values (\(1.0000, 0.7646, 0.7463, 0.7268, 0.5558, 0.5424, 0.4147\))
{[}2{]}.

\begin{center}\rule{0.5\linewidth}{0.5pt}\end{center}

\subsection{\texorpdfstring{3. Theorem Q2: Spontaneous Emergence of a
Single Complex Structure
\(\mathbb{C}\)}{3. Theorem Q2: Spontaneous Emergence of a Single Complex Structure \textbackslash mathbb\{C\}}}\label{theorem-q2-spontaneous-emergence-of-a-single-complex-structure-mathbbc}

A defining postulate of quantum mechanics is that state spaces are
Hilbert spaces over \(\mathbb{C}\) (not \(\mathbb{R}\)) and composite
systems combine via \(\otimes_{\mathbb{C}}\). Starting from a purely
real tensor \(T \in \mathbb{R}^{64\times 64\times 64}\), this structure
emerges in two stages:

\textbf{Theorem Q2 (Single Complex Structure Lock)}: 1. \textbf{Local
Complexification (Layer 1)}: Decomposing the \(2\times 2\) real inner
blocks (\(b \in \{0,1\},\; a,c \in \{4,5\}\)) of \(M_8\) in the Pauli
basis \(B = c_0 I + c_J J + c_x \sigma_x + c_z \sigma_z\)
(\(J = -i\sigma_y,\; J^2 = -I\)) yields \(|c_x|, |c_z| \sim 10^{-16}\)
(machine zero), leaving a pure complex field
\(\mathbb{C} \cong \{c_0 I + c_J J\}\) (Table 1). 2. \textbf{Global
Single-Clock Lock (Layer 2 \(\to\) Layer 3)}: A real tensor product of
two independent complex structures (Layer 2:
\(\mathbb{C}_{\mathrm{out}} \otimes_{\mathbb{R}} \mathbb{C}_{\mathrm{in}} \cong \mathbb{C}_+ \oplus \mathbb{C}_-\))
is unstable and collapses into Layer 3,
\[U_{\text{Layer 3}} = \frac{1}{\sqrt{2}}\Bigl(u \otimes v + (Ju)\otimes(-Jv)\Bigr),\]
locking \(J_{\mathrm{out}} = -J_{\mathrm{in}}\) into a single global
imaginary unit \(J\) (a single clock) while strictly doubling \(\|P\|\)
(Case 1 in {[}2{]}: \(0.240854 \to 0.481707\); verified here with
\(|\sigma_{1,2} - 1/\sqrt{2}| = 3.231 \times 10^{-7}\)).

{\def\LTcaptype{none} % do not increment counter
\begin{longtable}[]{@{}
  >{\raggedright\arraybackslash}p{(\linewidth - 4\tabcolsep) * \real{0.3333}}
  >{\raggedright\arraybackslash}p{(\linewidth - 4\tabcolsep) * \real{0.3333}}
  >{\raggedright\arraybackslash}p{(\linewidth - 4\tabcolsep) * \real{0.3333}}@{}}
\toprule\noalign{}
\begin{minipage}[b]{\linewidth}\raggedright
Verification Item (Theorem Q2)
\end{minipage} & \begin{minipage}[b]{\linewidth}\raggedright
Measured Value
\end{minipage} & \begin{minipage}[b]{\linewidth}\raggedright
Interpretation
\end{minipage} \\
\midrule\noalign{}
\endhead
\bottomrule\noalign{}
\endlastfoot
Layer 1 \(M_8\) (\(b=0\)) complex coeffs \(c_0, c_J\) &
\(+0.27343360,\; +0.49200681\) & Pure complex field \(\mathbb{C}\) \\
Layer 1 \(M_8\) (\(b=0\)) anisotropic \(|c_x|, |c_z|\) &
\(1.943 \times 10^{-16},\; 6.661 \times 10^{-16}\) & Exact zero (machine
precision) \\
Layer 1 \(M_8\) (\(b=1\)) anisotropic \(|c_x|, |c_z|\) &
\(5.551 \times 10^{-17},\; 5.274 \times 10^{-16}\) & Exact zero (machine
precision) \\
Layer 3 single-complex lock \(|\sigma_{1,2} - 1/\sqrt{2}|\) &
\(3.231 \times 10^{-7}\) (\(\|P\| \times 2\)) &
\(\mathbb{C}_{\mathrm{out}} \otimes_{\mathbb{R}} \mathbb{C}_{\mathrm{in}} \to \mathbb{C}_{\mathrm{single}}\) \\
\end{longtable}
}

\begin{center}\rule{0.5\linewidth}{0.5pt}\end{center}

\subsection{\texorpdfstring{4. Theorem Q3: \(O(N)\) Covariance, \(L^2\)
Uniqueness, and the Born
Rule}{4. Theorem Q3: O(N) Covariance, L\^{}2 Uniqueness, and the Born Rule}}\label{theorem-q3-on-covariance-l2-uniqueness-and-the-born-rule}

\textbf{Theorem Q3 (\(L^2\) Uniqueness and Born Rule \(P = |\psi|^2\))}:
For any orthogonal transformation \(R \in O(N)\) (\(R^T R = I\)), the
index contraction of \(\mathrm{Tri}\) satisfies strict covariance:
\[\mathrm{Tri}\bigl((R \otimes R \otimes R) U\bigr) = (R \otimes R \otimes R)\,\mathrm{Tri}(U).\]
Applying a deterministic Givens-product rotation \(R_{64} \in SO(64)\)
(\(\|R_{64}^T R_{64} - I\| = 1.890 \times 10^{-15}\)) to \(H\) confirms
a relative covariance error of \textbf{\(8.843 \times 10^{-16}\)} (Table
2, Fig. 1(b)).

Under \(O(N)\) covariance, the \textbf{unique} basis-invariant \(L^p\)
norm is \textbf{\(p = 2\)} (Hilbert/Euclidean norm). Under \(R_{64}\),
the \(L^1\) norm changes by \(+351.8\%\) and \(L^4\) by \(-42.0\%\),
whereas \(L^2\) is invariant to \textbf{\(2.220 \times 10^{-16}\)}. Thus
any basis-independent partitioning across orthogonal subspaces
(channels, sites, harmonics, spins, or Bell pairs) is fixed to the
squared-amplitude measure \(|c_m|^2\). Combined with the emergent
complex structure \(\mathbb{C}\) (Theorem Q2) and dimension \(d \ge 3\),
\textbf{Gleason's theorem} {[}4{]} uniquely forces every non-contextual
probability measure on orthogonal projections \(\Pi\) to take the
Born-rule form:
\[P(\Pi) = \mathrm{Tr}(\rho\,\Pi) = |\langle \phi | \psi \rangle|^2.\]
This directly explains why the spin-\(s\) amplitude law \(1/\sqrt{s+1}\)
(squared weights \(1:1/2:1/3\)) {[}3{]}, Bell-pair singular values
\((1/2)^2 \times 4 = 1\) {[}2{]}, and two-particle coexistence energies
{[}3{]} all obey exact \(L^2\) conservation.

{\def\LTcaptype{none} % do not increment counter
\begin{longtable}[]{@{}
  >{\raggedright\arraybackslash}p{(\linewidth - 6\tabcolsep) * \real{0.2105}}
  >{\centering\arraybackslash}p{(\linewidth - 6\tabcolsep) * \real{0.2632}}
  >{\centering\arraybackslash}p{(\linewidth - 6\tabcolsep) * \real{0.2632}}
  >{\centering\arraybackslash}p{(\linewidth - 6\tabcolsep) * \real{0.2632}}@{}}
\toprule\noalign{}
\begin{minipage}[b]{\linewidth}\raggedright
\(O(64)\) Rotation Test (Theorem Q3)
\end{minipage} & \begin{minipage}[b]{\linewidth}\centering
Before Rotation
\end{minipage} & \begin{minipage}[b]{\linewidth}\centering
After Rotation
\end{minipage} & \begin{minipage}[b]{\linewidth}\centering
Relative Change
\end{minipage} \\
\midrule\noalign{}
\endhead
\bottomrule\noalign{}
\endlastfoot
\(\mathrm{Tri}\) covariance error
\(\|\mathrm{Tri}(R^3 H)-R^3\mathrm{Tri}(H)\|/\|\mathrm{Tri}(H)\|\) & ---
& --- & \textbf{\(8.843 \times 10^{-16}\)} \\
\(L^1\) norm (\(p = 1\)) & \(21.531074\) & \(97.269545\) & \(+351.8\%\)
(non-invariant) \\
\textbf{\(L^2\) norm (\(p = 2\): Hilbert / Born)} &
\textbf{\(1.000000\)} & \textbf{\(1.000000\)} &
\textbf{\(2.220 \times 10^{-16}\) (exact invariant)} \\
\(L^4\) norm (\(p = 4\)) & \(0.227654\) & \(0.132001\) & \(-42.0\%\)
(non-invariant) \\
\end{longtable}
}

\begin{figure}[H]
\centering
\includegraphics[width=0.96\linewidth,height=\textheight,keepaspectratio,alt={Verification of Theorems Q2 and Q3. (a) Pauli decomposition of inner M\_8 in Layer 1 and single-complex-structure locking in Layer 3. (b) Exact O(64) covariance (8.84 \textbackslash times 10\^{}\{-16\}) and rotational invariance test of L\^{}p norms (p=1,2,4), uniquely singling out p=2.}]{figs_en/fig_r4_complex_and_born_en.png}
\caption{Verification of Theorems Q2 and Q3. (a) Pauli decomposition of
inner \(M_8\) in Layer 1 and single-complex-structure locking in Layer
3. (b) Exact \(O(64)\) covariance (\(8.84 \times 10^{-16}\)) and
rotational invariance test of \(L^p\) norms (\(p=1,2,4\)), uniquely
singling out \(p=2\).}
\end{figure}

\begin{center}\rule{0.5\linewidth}{0.5pt}\end{center}

\subsection{\texorpdfstring{5. Theorem Q4: \(\mathbb{Z}_2^3\) Discrete
Uncertainty Saturation and Minimum
Resolution}{5. Theorem Q4: \textbackslash mathbb\{Z\}\_2\^{}3 Discrete Uncertainty Saturation and Minimum Resolution}}\label{theorem-q4-mathbbz_23-discrete-uncertainty-saturation-and-minimum-resolution}

\textbf{Theorem Q4 (Donoho--Stark and Entropic Uncertainty Bounds)}: For
any normalized distribution \(f: \mathbb{Z}_2^3 \to \mathbb{R}\)
(\(\sum_x |f(x)|^2 = 1\)) on the outer crystal
\(O_8 \cong \mathbb{Z}_2^3\) and its normalized Walsh--Hadamard spatial
harmonic transform \(\hat{f} = H_8 f\), two exact uncertainty
inequalities hold {[}5{]}:
\[|\mathrm{supp}(f)| \cdot |\mathrm{supp}(\hat{f})| \;\ge\; 8, \qquad H_{\mathrm{space}}(f) + H_{\mathrm{harmonic}}(\hat{f}) \;\ge\; \ln 8 = 3\ln 2 \approx 2.079442.\]

As shown in Table 3 and Fig. 2, all three bare localization prototypes
identified in Rapid Communication 3 {[}3{]} \textbf{strictly saturate
both uncertainty bounds} (zero excess). By contrast, when evolved on the
time-crystal background \(H\) to \(\tau = 30\), the dressed stable
excitations (\(L_x\), P5, P4) strictly respect the bound while shifting
upward (\(+0.0512\) to \(+0.3130\)) in direct proportion to their
interaction clouds.

{\def\LTcaptype{none} % do not increment counter
\begin{longtable}[]{@{}
  >{\raggedright\arraybackslash}p{(\linewidth - 10\tabcolsep) * \real{0.1379}}
  >{\centering\arraybackslash}p{(\linewidth - 10\tabcolsep) * \real{0.1724}}
  >{\centering\arraybackslash}p{(\linewidth - 10\tabcolsep) * \real{0.1724}}
  >{\centering\arraybackslash}p{(\linewidth - 10\tabcolsep) * \real{0.1724}}
  >{\centering\arraybackslash}p{(\linewidth - 10\tabcolsep) * \real{0.1724}}
  >{\centering\arraybackslash}p{(\linewidth - 10\tabcolsep) * \real{0.1724}}@{}}
\toprule\noalign{}
\begin{minipage}[b]{\linewidth}\raggedright
State (on \(\mathbb{Z}_2^3\))
\end{minipage} & \begin{minipage}[b]{\linewidth}\centering
\(\|\mathrm{supp}(x)\| \times \|\mathrm{supp}(k)\|\)
\end{minipage} & \begin{minipage}[b]{\linewidth}\centering
\(H_{\mathrm{space}}\)
\end{minipage} & \begin{minipage}[b]{\linewidth}\centering
\(H_{\mathrm{harmonic}}\)
\end{minipage} & \begin{minipage}[b]{\linewidth}\centering
\(H_x + H_k\)
\end{minipage} & \begin{minipage}[b]{\linewidth}\centering
Excess over \(\ln 8\)
\end{minipage} \\
\midrule\noalign{}
\endhead
\bottomrule\noalign{}
\endlastfoot
\textbf{{[}Bare{]} Nonlocal wave (\(H\), P3a core)} &
\textbf{\(8 \times 1 = 8\)} & \(2.079442\) & \(0.000000\) &
\textbf{\(2.079442\)} & \textbf{\(0.000000\) (exact saturation)} \\
\textbf{{[}Bare{]} Quaternion core (\(H_4\))} &
\textbf{\(4 \times 2 = 8\)} & \(1.386294\) & \(0.693147\) &
\textbf{\(2.079442\)} & \textbf{\(0.000000\) (exact saturation)} \\
\textbf{{[}Bare{]} 1-site core (\(\delta_{000}\))} &
\textbf{\(1 \times 8 = 8\)} & \(0.000000\) & \(2.079442\) &
\textbf{\(2.079442\)} & \textbf{\(0.000000\) (exact saturation)} \\
{[}Dressed \(\tau=30\){]} Background \(H\) / P3a & \(8 \times 1 = 8\) &
\(2.079442\) & \(0.000000\) & \(2.079442\) & \(+0.000000\) (harm.
\(0.00\%\)) \\
{[}Dressed \(\tau=30\){]} Spin-1 gauge \(L_x\) & --- & \(2.058650\) &
\(0.072019\) & \(2.130669\) & \(+0.051228\) (harm. \(3.66\%\)) \\
{[}Dressed \(\tau=30\){]} \(H_4\) + cloud P5 & --- & \(2.022809\) &
\(0.129825\) & \(2.152634\) & \(+0.073192\) (harm. \(10.00\%\)) \\
{[}Dressed \(\tau=30\){]} 1-site + cloud P4 & --- & \(1.793664\) &
\(0.598750\) & \(2.392414\) & \(+0.312972\) (harm. \(45.85\%\)) \\
\end{longtable}
}

\textbf{Dimensionless Planck-Scale Prototypes}: The finite crystal
enforces a minimum spatial resolution \(\Delta x_{\min} = 1\) site
(Planck length \(\ell_P\) prototype), a maximum harmonic order
\(k_{\max} = 3\) (8-mode UV cutoff), a minimum update step
\(\Delta\tau_{\min} = 1\) (Planck time \(t_P\) prototype), and a minimum
phase-space cell \(\ln 8\) (dimensionless \(\hbar\) prototype),
eliminating continuum UV divergences at the axiomatic level.

\begin{figure}[H]
\centering
\includegraphics[width=0.96\linewidth,height=\textheight,keepaspectratio,alt={Verification of Theorem Q4. (a) \textbackslash mathbb\{Z\}\_2\^{}3 entropic uncertainty plane (H\_x + H\_k \textbackslash ge \textbackslash ln 8): bare cores strictly saturate the boundary, while \textbackslash tau=30 dressed states shift above the bound by their cloud weights. (b) Donoho--Stark support-product saturation xy = 8 with minimum resolution \textbackslash Delta x\_\{\textbackslash min\}=1 and UV cutoff k\_\{\textbackslash max\}=3.}]{figs_en/fig_r4_uncertainty_and_cutoff_en.png}
\caption{Verification of Theorem Q4. (a) \(\mathbb{Z}_2^3\) entropic
uncertainty plane (\(H_x + H_k \ge \ln 8\)): bare cores strictly
saturate the boundary, while \(\tau=30\) dressed states shift above the
bound by their cloud weights. (b) Donoho--Stark support-product
saturation \(xy = 8\) with minimum resolution \(\Delta x_{\min}=1\) and
UV cutoff \(k_{\max}=3\).}
\end{figure}

\begin{center}\rule{0.5\linewidth}{0.5pt}\end{center}

\subsection{6. Discussion, Symmetry Inventory, and
Outlook}\label{discussion-symmetry-inventory-and-outlook}

Theorems Q1--Q4 establish that the real 3-tensor triangle-loop product
intrinsically generates discretization, a single complex Hilbert
structure, the Born rule, and discrete uncertainty saturation with a UV
cutoff. To maintain strict scientific precision, we separate established
results from open problems:

\begin{itemize}
\tightlist
\item
  \textbf{Established (Rapid Reports 1--4)}: Hurwitz crystal selection
  (\(O_8\)), 4-stage doubling-entanglement cascade and Bell pairs,
  \(SO(3)\) spin selection (\(1\) vs.~\((1/3)^\tau\)) and
  \(1/\sqrt{s+1}\) law, \(G_2\) (\(14\)) \(\supset SU(3)\) (\(8\)
  gluons) and \(U(1)\) phase symmetries, two-particle coexistence, and
  the four quantum structure theorems Q1--Q4.
\item
  \textbf{Open Problems for Subsequent Reports}: (1) Quantitative
  relation between localized energy, local time dilation (first seen as
  P4 slowdown \texttt{d1:\ 0.904\ -\textgreater{}\ 0.120} in {[}2{]}),
  signal-propagation speed limit \(c\), and Lorentz/gravitational
  factors; (2) explicit chiral \(SU(2)_L\) parity violation; and (3)
  fermionic spin-\(1/2\) exchange antisymmetry via octonionic triality
  (\(8_v, 8_s, 8_c\)).
\end{itemize}

\begin{center}\rule{0.5\linewidth}{0.5pt}\end{center}

\subsection{7. Reproducibility and Released
Files}\label{reproducibility-and-released-files}

All numerical results and figures are deterministically reproducible
without randomness: - Verification script:
\texttt{tri\_loop\_n64\_quantum\_foundations\_verify\_v0.py}
(\textasciitilde15 s) - Numerical log:
\texttt{results/tri\_loop\_n64\_quantum\_foundations\_verify\_v0.txt} -
Figure generator: \texttt{plot\_fourth\_report\_figures\_v0.py} -
Figures: \texttt{results/fig\_r4\_complex\_and\_born\_\{ja,en\}.png} and
\texttt{results/fig\_r4\_uncertainty\_and\_cutoff\_\{ja,en\}.png} - TeX
builder: \texttt{build\_tex\_rapid4.py}

\subsection{References}\label{references}

\begin{enumerate}
\def\labelenumi{\arabic{enumi}.}
\tightlist
\item
  N. Kihara, \emph{Self-Consistent Crystal Selection, Dimensional
  Melting, and Two-Scale Particle Condensation in the Triangle-Loop
  Product of 3D Cubic Arrays (Rapid Communication)}, Zenodo (2026), DOI:
  10.5281/zenodo.24616931.
\item
  N. Kihara, \emph{Four-Stage Doubling-Entanglement Phase Transitions,
  Single-Clock Complex Locking, Exact Triple-Outer-Product Bell Pairs,
  and Probe-Type Branching in the 3D Triangle-Loop Product (Second Rapid
  Communication)}, Zenodo (2026), DOI: 10.5281/zenodo.24618400.
\item
  N. Kihara, \emph{SO(3) Spin Selection, G2/SU(3) Gauge Symmetries,
  Spatial Harmonic Localization, and Two-Particle Coexistence in the 3D
  Triangle-Loop Product (Third Rapid Communication)}, Zenodo (2026),
  DOI: 10.5281/zenodo.24619357.
\item
  A. M. Gleason, \emph{Measures on the Closed Subspaces of a Hilbert
  Space}, J. Math. Mech. \textbf{6}, 885--893 (1957).
\item
  D. L. Donoho and P. B. Stark, \emph{Uncertainty Principles and Signal
  Recovery}, SIAM J. Appl. Math. \textbf{49}, 906--931 (1989).
\end{enumerate}

\end{document}
